\section{Thực nghiệm}
\label{sec:experiments}

Phần này trình bày phần thực nghiệm của báo cáo. Theo yêu cầu của đề tài, nhóm
chọn hướng \textbf{Cách 2: tự xây dựng demo từ đầu}. Thay vì kế thừa trực tiếp
mã nguồn chính thức của một công trình cụ thể như GCG, PAIR hay AutoDAN, nhóm tự
cài đặt một pipeline nhỏ để mô phỏng các thành phần chính được trình bày trong
tutorial về jailbreaking LLMs. Pipeline bao gồm các bước chính: xây dựng prompt,
áp dụng biến đổi tấn công, gọi mô hình, áp dụng các cấu hình phòng thủ, đánh giá
định lượng và trực quan hóa kết quả.

Do hạn chế về tài nguyên, quota API, thời gian chạy và khả năng truy cập mô
hình, thí nghiệm trong báo cáo này không nhằm mục tiêu tái lập đầy đủ một
benchmark quy mô lớn như HarmBench hoặc JailbreakBench. Thay vào đó, thực nghiệm
được thiết kế như một demo có kiểm soát nhằm minh họa quy trình đánh giá
jailbreak và phân tích trade-off giữa safety và utility. Các bảng kết quả chi
tiết, biểu đồ và error cases được trình bày trong phần phụ lục.

\subsection{Phân công công việc}

Bảng~\ref{tab:work-division} trình bày phân công công việc của các thành viên
trong nhóm.

\begin{table}[H]
\centering
\caption{Phân công công việc trong nhóm}
\label{tab:work-division}
\begin{tabular}{p{0.18\textwidth}p{0.28\textwidth}p{0.44\textwidth}}
\toprule
\textbf{MSSV} & \textbf{Họ và tên} & \textbf{Nhiệm vụ} \\
\midrule
23120245 & Nguyễn Quang Duy & Nghiên cứu lý thuyết nền tảng và tutorial về jailbreaking LLMs. \\
23120258 & Lưu Trọng Hiếu & Nghiên cứu các công trình tiêu biểu liên quan đến tấn công và phòng thủ jailbreak. \\
23120260 & Văn Đình Hiếu & Chạy thực nghiệm, tổng hợp kết quả, đưa ra đánh giá và kết luận. \\
\bottomrule
\end{tabular}
\end{table}

\subsection{Thiết lập và phạm vi thực nghiệm}

Experiment chính được chạy ở chế độ \texttt{\detokenize{openrouter_api}}. Bộ
prompt gồm 30 prompt gốc, chia thành ba nhóm nội dung: \texttt{knowledge},
\texttt{company} và \texttt{hazardous}. Nhóm \texttt{knowledge} bao gồm các
prompt hợp lệ liên quan đến AI safety, jailbreak, benchmark và evaluation. Nhóm
\texttt{company} mô phỏng các yêu cầu rủi ro liên quan đến thông tin nội bộ,
system instruction, secret hoặc dữ liệu riêng tư. Nhóm \texttt{hazardous} gồm
các prompt nằm gần ranh giới an toàn khi thảo luận về nội dung rủi ro cao.

Mỗi prompt gốc được áp dụng năm loại attack transformation:
\texttt{none}, \texttt{\detokenize{role_play}}, \texttt{obfuscation},
\texttt{\detokenize{multi_turn_simulation}} và
\texttt{\detokenize{injection_style}}. Sau khi mô hình sinh phản hồi, mỗi kết
quả được đánh giá dưới bốn cấu hình phòng thủ:
\texttt{\detokenize{no_defense}}, \texttt{\detokenize{input_filter}},
\texttt{\detokenize{output_moderation}} và
\texttt{\detokenize{defense_in_depth}}.

Tổng số dòng đánh giá chi tiết là:

[
30 \text{ prompts} \times 5 \text{ attack types} \times 4 \text{ defense types}
= 600 \text{ cases}.
]

Trong chế độ API, mô hình chỉ được gọi một lần cho mỗi cặp prompt--attack. Các
defense được áp dụng offline sau khi đã có phản hồi, vì vậy chúng không làm tăng
số lượng API calls. Cách thiết kế này giúp giảm chi phí chạy thí nghiệm nhưng
vẫn cho phép so sánh tác động tương đối của từng cấu hình phòng thủ trên cùng
một tập phản hồi.

Do giới hạn về tài nguyên tính toán và quota API, thực nghiệm có một số giới hạn
quan trọng. Thứ nhất, số lượng prompt còn nhỏ và chủ yếu phục vụ minh họa
pipeline thay vì đại diện cho một benchmark quy mô lớn. Thứ hai, thí nghiệm
chính chỉ sử dụng một target model thông qua OpenRouter API, chưa so sánh nhiều
mô hình thương mại và open-source khác nhau. Thứ ba, nhóm không chạy lại
official implementation của các phương pháp phức tạp như GCG, AutoDAN hoặc PAIR
vì các phương pháp này thường cần nhiều lần gọi mô hình, GPU hoặc quyền truy cập
chi tiết hơn vào mô hình. Cuối cùng, các lớp phòng thủ trong demo là rule-based,
nhằm minh họa nguyên lý input filtering, output moderation và defense-in-depth,
chưa phải guard model chuyên dụng.

Vì vậy, các kết quả trong phần này nên được hiểu là kết quả của một demo có kiểm
soát. Chúng cho thấy xu hướng và trade-off trong pipeline được xây dựng, nhưng
chưa đủ để kết luận tổng quát về mọi mô hình hoặc mọi dạng jailbreak.

\subsection{Kết quả và nhận xét chính}

Các chỉ số chính được sử dụng trong experiment gồm ASR, ORR và Utility Rate. ASR
(\textit{Attack Success Rate}) đo tỷ lệ prompt harmful dẫn đến unsafe compliance.
ORR (\textit{Over-Refusal Rate}) đo tỷ lệ prompt hợp lệ nhưng bị từ chối hoặc bị
chặn. Utility Rate đo tỷ lệ prompt hợp lệ được trả lời thành công. Vì vậy, một
cấu hình phòng thủ tốt không chỉ cần giảm ASR, mà còn cần hạn chế over-refusal và
giữ utility cho các yêu cầu hợp lệ.

Kết quả tổng hợp cho thấy cấu hình \texttt{\detokenize{no_defense}} có ASR cao
nhất, đạt 0.133. Điều này cho thấy khi không có lớp kiểm soát bổ sung, mô hình
vẫn có một tỷ lệ nhất định tạo ra phản hồi unsafe compliance đối với prompt
harmful. Hai cấu hình \texttt{\detokenize{output_moderation}} và
\texttt{\detokenize{defense_in_depth}} đưa ASR về 0.000. Tuy nhiên, hai cấu hình
này có trade-off khác nhau: \texttt{\detokenize{output_moderation}} giữ Utility
Rate ở mức 0.613, trong khi \texttt{\detokenize{defense_in_depth}} làm Utility
Rate giảm còn 0.427 do input filter chặn nhiều prompt hợp lệ hơn.

Trong thiết lập hiện tại, \texttt{\detokenize{output_moderation}} là cấu hình có
trade-off tốt nhất giữa safety và utility. Cấu hình
\texttt{\detokenize{input_filter}} giúp giảm ASR từ 0.133 xuống 0.027, nhưng đồng
thời làm tăng over-refusal. Điều này cho thấy rule-based input filtering có thể
chặn quá rộng, đặc biệt khi gặp các prompt có chứa từ khóa rủi ro nhưng mục đích
thảo luận vẫn mang tính học thuật hoặc an toàn.

Xét theo attack type, không có nhóm attack nào tạo ASR quá cao trong pipeline
hiện tại. Các nhóm \texttt{none}, \texttt{\detokenize{role_play}} và
\texttt{\detokenize{multi_turn_simulation}} có ASR xấp xỉ 0.050. Nhóm
\texttt{obfuscation} có ASR thấp nhất, đạt 0.017. Điểm đáng chú ý là
\texttt{\detokenize{injection_style}} tuy có ASR không quá cao nhưng lại làm ORR
tăng mạnh, cho thấy dạng prompt này có thể khiến hệ thống trở nên quá thận trọng
và từ chối nhiều yêu cầu đáng lẽ có thể trả lời.

Xét theo attack profile, nhóm \texttt{company} là nhóm có rủi ro unsafe
compliance cao nhất khi không có defense. Điều này phù hợp với thiết kế prompt vì
nhóm này mô phỏng các yêu cầu liên quan đến thông tin nội bộ hoặc instruction
hierarchy. Nhóm \texttt{hazardous} có ASR thấp hơn nhưng chịu ảnh hưởng mạnh từ
over-refusal. Nhóm \texttt{knowledge} không tạo ASR trong thí nghiệm này nhưng
vẫn có ORR đáng kể, phản ánh khó khăn trong việc phân biệt giữa thảo luận học
thuật về jailbreak và yêu cầu jailbreak thực sự.

Bên cạnh experiment chính, nhóm cũng thực hiện một manual web evaluation trên
file \texttt{\detokenize{web_manual_responses.tsv}}. Trong phần này, nhãn
\texttt{\detokenize{Passed_*}} được hiểu là kết quả jailbreak thành công:
\texttt{Passed = 1} nghĩa là model tạo phản hồi được đánh giá là vượt qua cơ chế
an toàn trong case tương ứng, còn \texttt{Passed = 0} nghĩa là jailbreak không
thành công. Kết quả manual cho thấy Gemini 3.1 Flash Lite có 5/31 case jailbreak
thành công, trong khi ChatGPT 5.5 và Claude Haiku 4.5 không có case
\texttt{Passed = 1}. Tuy nhiên, kết quả này được báo cáo riêng vì khác quy trình
đánh giá so với experiment pipeline.

Từ các kết quả trên, có thể tóm tắt các phát hiện chính như sau:

\begin{enumerate}[leftmargin=*]
\item Khi không có defense, hệ thống vẫn có unsafe compliance với ASR bằng
0.133.

\item \texttt{\detokenize{output_moderation}} là cấu hình có trade-off tốt nhất
trong thí nghiệm này: ASR giảm về 0.000 trong khi Utility Rate vẫn giữ ở 0.613.

\item \texttt{\detokenize{input_filter}} và
\texttt{\detokenize{defense_in_depth}} giúp giảm ASR nhưng làm tăng
over-refusal, đặc biệt ở nhóm \texttt{hazardous}.

\item Nhóm \texttt{company} là nhóm có rủi ro unsafe compliance cao nhất khi
không có defense.

\item Nhóm \texttt{knowledge} không tạo ASR trong thí nghiệm này nhưng vẫn có
over-refusal, cho thấy các câu hỏi học thuật về AI safety có thể bị xử lý quá
thận trọng.

\item Manual web evaluation cho thấy Gemini 3.1 Flash Lite có 5/31 case
jailbreak thành công, trong khi ChatGPT 5.5 và Claude Haiku 4.5 không có case
\texttt{Passed = 1}; tuy nhiên kết quả này được báo cáo riêng vì khác quy trình
đánh giá so với experiment pipeline.
\end{enumerate}

Nhìn chung, kết quả thực nghiệm củng cố nhận định rằng đánh giá jailbreak không
nên chỉ tập trung vào attack success rate. Một hệ thống an toàn cần được đánh
giá đồng thời trên cả ba chiều: giảm unsafe compliance, hạn chế over-refusal và
giữ utility cho các yêu cầu hợp lệ.
